\begin{apartats}

\apartat{0,75}
Determina el domini i estudia la continuïtat de les funcions següents:
\begin{graella}{2}
  \sa y=\sqrt{x^2-2x-3} & \sa y=\ln|x-2|
\end{graella}

\begin{solucio}
i) Cal $x^2-2x-3=(x-3)(x+1)\ge0$: $\mathrm{Dom}=(-\infty,-1]\cup[3,+\infty)$.
És contínua a tot el domini.\\
ii) Cal $|x-2|>0$, és a dir $x\ne2$: $\mathrm{Dom}=\mathbb{R}\setminus\{2\}$.
És contínua a tot el domini. En $x=2$ els dos laterals valen $-\infty$:
\textbf{salt infinit} (asímptota vertical $x=2$).
\end{solucio}

\apartat{1}
Troba els punts de discontinuïtat de la funció
\[
f(x)=\frac{x-1}{x^3+3x^2-4}
\]
i classifica'ls.

\begin{solucio}
$x=1$ anul·la el denominador. Per Ruffini, $x^3+3x^2-4=(x-1)(x^2+4x+4)=(x-1)(x+2)^2$.
El domini és $\mathbb{R}\setminus\{-2,1\}$ i, per a $x\neq1$, $f(x)=\dfrac{1}{(x+2)^2}$.\\
$x=1$: $\lim_{x\to1}f(x)=\dfrac19$, però $f(1)$ no existeix: \textbf{evitable}.\\
$x=-2$: $(x+2)^2>0$ als dos costats, així que els dos laterals valen $+\infty$:
\textbf{salt infinit}, amb $\lim_{x\to-2}f(x)=+\infty$.
\end{solucio}

\apartat{0,75}
\textbf{Inventa.} Escriu una funció definida a trossos que sigui contínua a tot
$\mathbb{R}$ excepte en $x=0$, on ha de tenir una discontinuïtat evitable, i en $x=3$, on
ha de tenir una discontinuïtat de salt finit.

\begin{solucio}
Per exemple:
\[
f(x)=\begin{cases} x & \si{x<3,\ x\ne0},\\ 1 & \si{x=0},\\ x+1 & \si{x\ge3}.\end{cases}
\]
En $x=0$: $\lim_{x\to0}f(x)=0\ne f(0)=1$, evitable. En $x=3$: laterals $3$ i $4$, salt
finit. A la resta de punts, les branques són polinòmiques. (La resposta no és única.)
\end{solucio}

\end{apartats}
